\documentclass[10pt,a4paper]{amsart}
\usepackage[margin=19mm]{geometry}
\usepackage{amsmath,amssymb,mathtools}
\usepackage[scr=rsfs,cal=euler]{mathalfa}
\usepackage{xfrac}
\usepackage[unicode,hidelinks,bookmarks=false]{hyperref}
\newcommand{\ie}{i.e.\ }
\newcommand{\cf}{cf.\ }
\newcommand{\resp}{respectively\ }
\newcommand{\Subsection}{\subsection}
\providecommand{\nobreakdash}{\nobreak\hbox{-}}
\setlength{\emergencystretch}{2em}
\multlinegap0pt
\usepackage{amssymb}
\usepackage{bm}
\usepackage{xcolor}
\usepackage[shortlabels]{enumitem}
\newtheorem{proposition}{Proposition}
\DeclareMathOperator{\curl}{curl}
\DeclareMathOperator{\dive}{div}
\title{full title}
\author{Qian Huang, Christian Rohde, Ruixi Zhang}
\date{Source: arXiv:2601.19846v1 (CC BY 4.0)}
\begin{document}
\maketitle

\noindent\emph{Source and licence.} full title. Version arXiv:2601.19846v1 (CC BY 4.0), distributed by arXiv under the Creative Commons Attribution 4.0 International licence, http://creativecommons.org/licenses/by/4.0/, as recorded in the arXiv licence metadata for this exact version and shown on the abstract page https://arxiv.org/abs/2601.19846v1. Full licence notice: \texttt{licenses/arxiv-2601.19846-CC-BY-4.0.txt}.\par\smallskip

\noindent\textit{Editorial scope.} Counted unit: the article's proposition eq:F0\_evolve (source0.tex lines 837--851). The article's complete original proof of it is delivered with it (source0.tex lines 853--872, one \texttt{proof} environment of 20 lines). No mathematics is written, repaired or re-derived: the statement and the proof are the article's own lines, in the article's own notation, and no retained line is substituted or re-flowed.  Retained uncounted as the source-local context the proof needs: lines 811--813; lines 815--817; lines 819--821.  Nothing was omitted from any retained span, so the package carries no omission record.  No retained line contains a citation, so the wrapper carries no bibliography and no bibliography entry.  No retained line contains a figure environment, an image-inclusion command or drawing code, so no image is delivered and the package declares no asset dependency.  Editorial formatting only, in addition: an AMS \texttt{amsart} preamble that restates the article's own declarations and macros used by the retained lines, namely the environments \texttt{proposition} on separate counters, together with the standard packages those lines need.  The retained environments print in this document as Proposition 1 in order of appearance, whereas the article prints the same items with the numbering of its own full document.  The complete native source of the article as distributed in the arXiv e-print for this exact version is bundled unchanged as \texttt{sources/193491/source0.tex}.
\begin{equation} \label{eq:A0_def}
    A:=\curl\dive(u\otimes u) = u\cdot \nabla \omega - \omega \cdot \nabla u,
\end{equation}

\begin{equation} \label{eq:omega_ns}
    \partial_t\omega - \Delta \omega = -A.
\end{equation}

\begin{equation} \label{eq:omega_ed}
    \partial_t \left( \omega^{\epsilon,\delta} + \delta\partial_t\omega^{\epsilon,\delta} \right) - \Delta \omega^{\epsilon,\delta} =  A^{\epsilon,\delta}, \qquad A^{\epsilon,\delta}:= \curl\dive(u^{\epsilon,\delta}\otimes u^{\epsilon,\delta}).
\end{equation}

\begin{proposition}
The energy $\bar F^{\epsilon,\delta}$ and $F^{\epsilon,\delta}$ satisfy
\begin{equation} \label{eq:F0_evolve}
    \frac{d}{dt}\frac{1}{2}\bar F^{\epsilon,\delta}(t) = -\int_{\mathbb T^3} \left( \delta \left|\partial_t \omega^{\epsilon,\delta} \right|^2 + \left| \nabla\omega^{\epsilon,\delta} \right|^2 \right) \, dx -\int_{\mathbb T^3} \left( \omega^{\epsilon,\delta}+2\delta\partial_t \omega^{\epsilon,\delta} \right)\cdot A^{\epsilon,\delta} \, dx
\end{equation}
and
\begin{equation} \label{eq:modu_energy_evolve}
\begin{split}
    \frac{d}{dt}\frac{1}{2}F^{\epsilon,\delta}(t) &= T^{\epsilon,\delta}_\text{I}(t) + T^{\epsilon,\delta}_\text{II}(t) - \int_{\mathbb T^3}\left( \left| \nabla(\omega^{\epsilon,\delta}-\omega) \right|^2 + \delta \left| \partial_t\omega^{\epsilon,\delta} \right|^2 \right)\,dx, \\
    T^{\epsilon,\delta}_\text{I}(t) &:= \int_{\mathbb T^3} \delta\partial_t\omega^{\epsilon,\delta} \cdot\left(A-\Delta\omega-2A^{\epsilon,\delta} \right) \, dx, \\
    T^{\epsilon,\delta}_\text{II}(t) &:= - \int_{\mathbb T^3} (\omega^{\epsilon,\delta}-\omega)\cdot(A^{\epsilon,\delta}-A)\,dx,
\end{split}
\end{equation}
with  $A$ from \eqref{eq:A0_def} and $A^{\epsilon,\delta}$ from \eqref{eq:omega_ed}.
\end{proposition}

\begin{proof}
First, (\ref{eq:F0_evolve}) can be derived by multiplying (\ref{eq:omega_ed}) with $\omega^{\epsilon,\delta} + 2\delta \partial_t \omega^{\epsilon,\delta}$ and integrating (by parts) over $\mathbb T^3$.

Then, we note that
\[
    F^{\epsilon,\delta}(t) = \bar F^{\epsilon,\delta}(t) + \int_{\mathbb T^3} |\omega|^2\,dx - 2\int_{\mathbb T^3} \omega \cdot \left( \omega^{\epsilon,\delta} +\delta\partial_t \omega^{\epsilon,\delta} \right) \,dx.
\]
By using \eqref{eq:omega_ns} we have
\[
    \frac{d}{dt}\frac{1}{2}\int_{\mathbb T^3}|\omega|^2\,dx = -\int_{\mathbb T^3}|\nabla \omega|^2\,dx - \int_{\mathbb T^3}\omega\cdot A\, dx
\]
and hence
\[
\begin{aligned}
    \frac{d}{dt}\frac{1}{2}F^{\epsilon,\delta}(t) =& \frac{d}{dt}\frac{1}{2}\bar F^{\epsilon,\delta}(t) - \int_{\mathbb T^3}|\nabla \omega|^2\,dx - \int_{\mathbb T^3}\omega\cdot A\, dx \\
    &+ \int_{\mathbb T^3} \left( (A-\Delta\omega) \cdot \left(\omega^{\epsilon,\delta} + \delta \partial_t \omega^{\epsilon,\delta}\right) + \omega \cdot (A^{\epsilon,\delta}-\Delta \omega^{\epsilon,\delta}) \right) \,dx,
\end{aligned}
\]
where (\ref{eq:omega_ed}) and (\ref{eq:omega_ns}) were used. Then, (\ref{eq:modu_energy_evolve}) is derived by substituting (\ref{eq:F0_evolve}) into the above equation, integrating by parts, and rearranging the terms.
\end{proof}

\end{document}
